\chapter[Preservation, Discharge, and Transformation]{Preservation, Discharge, and the Transformation of Obligations}
\label{ch:preservation}

The preceding chapter showed that a successful verification leaves obligations that the kernel cannot discharge, and that for sufficiently expressive specification classes no total procedure can discharge them all. The question that follows is practical as well as theoretical. Verified artifacts do not stay as they are. They are translated, revised, repaired, specialized, combined with other artifacts, and imported into further developments. At each transformation the obligations attached to the input must be accounted for in the output. This chapter develops the central constructive notion of the monograph, which is the preservation of obligations through transformations, and separates it from the notion of discharge with which it is routinely confused. The formal definitions are those of Appendix~\ref{app:formal}, and the theorems proved there are restated here with their interpretation and their limits.

\section{Why traceability is not enough}

A natural first proposal for tracking obligations through a transformation is to require that every obligation of the input be accounted for in the output: either it is recorded as discharged, or it has a successor. This is traceability, and it is mechanically checkable. It is also insufficient, and the insufficiency is the reason the rest of the theory is needed.

Suppose an input obligation asserts a bound requiring $m+4$ derivatives of the data, and the transformation replaces it by a successor asserting the corresponding bound with $m+5$. The input obligation has a successor, and the ledger is complete. But the successor is a weaker claim. A reader who inspects the ledger and finds every obligation accounted for will conclude that nothing has been lost, and that conclusion is false. What the ledger lacks is a record that the content of the successor suffices for the content of the original. Traceability records where an obligation went. Preservation records whether what it became is enough.

Accordingly the theory has two layers. The syntactic layer concerns the form of the ledger and is decidable by inspection. The semantic layer concerns entailment between contents and is certified, not computed. The two layers have separate theorems, and keeping them separate is the safeguard against the error just described, namely mistaking a well-formed ledger for a semantically adequate one.

\section{Transitions, certificates, and the entailment condition}

A transformation $f_i$ from an artifact $X_{i-1}$ to an artifact $X_i$ is accompanied by a transition certificate $T_i = (Q_i, \Delta_i, \Theta_i)$. The relation $Q_i$ links obligations of the input to their successors in the output. The set $\Delta_i$ lists input obligations recorded as discharged at the transformation, each with a certificate that is currently warranted in the ledger state preceding it. The set $\Theta_i$ lists the assumptions the transformation is permitted to use. The certified context of an input obligation $o$ is the set consisting of the contents of its successors, the contents of the discharged obligations, and the permitted assumptions. The transformation \emph{preserves} $o$ when the certified context entails the content of $o$:
\[
\Ctx_i(o) \models_{\mathcal{I}} \varphi(o).
\]
It is preserving when it preserves every input obligation.

Two features of this definition deserve emphasis. First, the successors enter as premises along with the discharged obligations and the assumptions. Nothing requires that any one successor entail the original, and a conjunction of weaker successors may jointly entail it. The test is applied to the total certified context, which is the reading of RV-019 given in Section~\ref{sec:weakening17}. Second, the permitted assumptions are restricted. Without a restriction the condition is vacuous, because an assumption equal to $\varphi(o)$ would secure preservation of $o$ for free. Definition~\ref{def:admissible} of the appendix requires that each assumption either belong to the authorized context fixed with the instance or be the content of a tracked assumption obligation introduced by the transition, and that the dependency relation remain acyclic. An assumption introduced during a transformation therefore enters the residue and remains there until it is discharged. Preservation on the strength of an untracked assumption is not permitted, and preservation on the strength of a tracked one is permitted and honestly reported.

\section{The typing of transitions}

At each obligation the certified context stands in one of four relations to the original content, and the transition is typed accordingly. Write $\Ctx$ for the certified context of $o$. If $\Ctx \models \varphi(o)$ and $\varphi(o) \models \Ctx$, the transition is \emph{conservative}: what was required is exactly what is now certified. If $\Ctx \models \varphi(o)$ but $\varphi(o) \not\models \Ctx$, it is \emph{strengthening}: the original is preserved, and the context contains claims the original did not, which are a new burden to be discharged. If $\Ctx \not\models \varphi(o)$ but $\varphi(o) \models \Ctx$, it is \emph{weakening}: everything now certified follows from the original, and the original does not follow from it, so content has been lost. If neither entailment holds, it is \emph{lateral}: the transition has replaced the obligation by something incomparable. Conservative and strengthening transitions are preserving. Weakening and lateral transitions are not. The classification is exhaustive and exclusive by construction, and it is recorded at the level of individual obligations, with a transformation counted as preserving only if every obligation is handled by a preserving case.

\section{Traceability: the syntactic theorem}

A ledger is \emph{well formed} when every transition records a disposition for every inherited obligation, that is, each input obligation is either in $\Delta_i$ or has at least one successor. For an obligation $o$ of the initial artifact, the successor sets $S_k(o)$ are obtained by composing the successor relations. The following is Theorem~\ref{thm:trace} of the appendix.

\begin{theorem}[RV-020]
In a well-formed ledger, every originating obligation either has a certified discharge in its ancestry or has a nonempty set of successors at the final artifact. Well-formedness and these two alternatives are decidable for finite ledgers by inspection of the successor relations and discharge records.
\end{theorem}

\begin{corollary}[RV-017]
In a well-formed ledger, an originating obligation with no certified discharge in its ancestry has a successor at the final artifact. No obligation leaves the ledger without a recorded disposition. A discharge that is not recorded as a member of some $\Delta_k$ lies outside the formalism, and the corollary does not exclude it.
\end{corollary}

The proof is an induction on the length of the chain and is given in the appendix. The theorem is intentionally weak. It says that nothing disappears without record, and it says nothing about whether the recorded dispositions are adequate. Its value lies in being checkable even when the entailment relation is not. A Lean formalization of the ledger structure can verify well-formedness and compute the ancestry sets without any commitment about the semantics of the contents. The result is therefore the part of the theory that can be machine checked today. It is also the part most likely to be mistaken for the whole, which is why the next theorem is stated alongside it.

\section{Compositional preservation: the semantic theorem}

\begin{theorem}[RV-008]
\label{thm:rv008-ch}
Let $T_1, \dots, T_n$ be transition certificates over a consequence relation satisfying reflexivity, monotonicity, and cut, and for $o$ an obligation of the initial artifact let each $T_k$ be preserving at every obligation reached from $o$. let $\Gamma_n(o)$ be the set consisting of the contents of the obligations in $S_n(o)$ together with all discharged contents and permitted assumptions accumulated along the chain. Then $\Gamma_n(o) \models_{\mathcal{I}} \varphi(o)$.
\end{theorem}

The proof, given in the appendix as Theorem~\ref{thm:compose}, is an induction on $n$. In outline: at each step the context of any intermediate obligation is contained in the cumulative context, so by monotonicity the cumulative context entails the content of every intermediate obligation; the cumulative context of the previous step is thereby entailed element by element, and cut transfers the earlier entailment of $\varphi(o)$ to the new cumulative context.

The theorem is conditional, and the conditions are visible in its statement. It requires preservation at the obligations reached from $o$ and the three structural properties of the consequence relation. It does not require that the consequence relation be decidable, and it does not assert that any premise is true. This last point is what makes the entailment a result about preservation and not about truth. Admissibility of the certificates, soundness of the certificates, and soundness of the consequence relation for the interpretation enter only the second statement of the appendix theorem, which concerns establishment. The premises of the entailment are the contents of the composite successors $S_n(o)$, the contents of the obligations recorded in discharge sets along the chain, and the permitted assumptions. If all of these are established at the evaluation time and the soundness hypotheses hold, then $\varphi(o)$ is transported, in the appendix's sense that it holds under the interpretation. A member of $S_n(o)$ that is not warranted-discharged belongs to the residue of the final artifact. If the content of a recorded discharge or a tracked assumption is not established, its source obligation is not warranted-discharged and belongs to the transport shortfall, and to the residue of the final artifact only if the ledger carries it forward; a defeated authorized assumption has no source obligation and lies in no residue. In these cases $\varphi(o)$ has been preserved and has not been transported.

\section{Preservation is not discharge}

The independence of the two notions is the contribution that the monograph regards as central, and it is stated as a proposition rather than left as a slogan.

\begin{proposition}[RV-018]
\label{prop:rv018-ch}
Preservation of an obligation need not entail its discharge, and discharge of one obligation need not preserve another. An obligation recorded in the discharge set of a transition is preserved at that transition.
\end{proposition}

The appendix proves the first two parts by explicit ledgers. For the first, a chain in which each step replaces an open obligation by an identical open successor is preserving at every step by reflexivity and discharges nothing. For the second, a transition that discharges one obligation and leaves another with a successor too weak to recover it is well formed and non-preserving at the second, so discharge of the first supplies a premise and does not supply the content missing from the second. A third part records that an obligation recorded in the discharge set of a transition is trivially preserved at that step, because its content belongs to its own certified context.

The first half has a corollary that gives the monograph its title in miniature.

\begin{corollary}[RV-009]
A verified continuation can retain unresolved originating obligations.
\end{corollary}

A system can therefore be rigorous about an obligation it has not settled. It can carry a statement correspondence obligation through a sequence of translations and repairs, certify at each step that nothing has been lost, and arrive at the final artifact with the obligation exactly as open as it began. This is a precise sense in which a verification pipeline may be sound without being complete: it never reports as settled what it has only preserved. The value of such a pipeline is that the final report is honest about the residue, and a reader may decide whether the residue is acceptable.

\section{Weakening: the Navier-Stokes cases}
\label{sec:weakening17}

The preservation condition is a condition on the total certified context. A successor that is strictly weaker than its predecessor does not by itself make a transition non-preserving, because other successors, discharged obligations, or admissible assumptions may supply what the weaker successor lacks. Conversely a transition with a single weaker successor, no discharges, and no assumptions is non-preserving. This is Proposition~\ref{prop:weak} of the appendix, and it is the formal statement of RV-019.

The derivative case of Chapter~12 is the paradigm. Let $o$ be the obligation whose content is the natural-language estimate (8.19) of the source,
\[
\|N^{-1}F\|_{C^m_y} \le C_m \|F\|_{C^{m+4}_y},
\]
and let $o'$ be the obligation whose content is the Lean estimate
\[
\|N^{-1}F\|_{C^m_y} \le K_m \|F\|_{C^{m+5}_y},
\]
cited at the pinned commit as \texttt{norm\_derivativeWord\_inverse\_le}. Since $\|F\|_{C^{m+4}_y} \le \|F\|_{C^{m+5}_y}$, the first bound implies the second with $K_m = C_m$, so $\varphi(o) \models \varphi(o')$. The source reports the second as strictly weaker, that is, as not implying the first. With $\Delta = \Theta = \varnothing$ and $o'$ the only successor, the certified context is $\{\varphi(o')\}$ and the transition from the natural-language paper to the formalization is of weakening type at $o$. The classification depends on the failure of the converse entailment, which is the source's claim, and the monograph records it as such pending independent verification of that claim. The entailment relation here is derivability from the certified background of the development and not classical consequence over truths, as Chapter~\ref{ch:derivative} explains, and without that reading the classification would be vacuous. The classification also depends on the premise set. If the formalization contained a further certified lemma deriving the $m+4$ bound, the context would change and the type with it. The type is a statement about the whole context, and the example is accordingly evaluated against the full set of premises available in the formalized development and not against one declaration in isolation.

The pressure-flux case of Chapter~13 is a different type. The natural-language estimate (10.19) bounds the pressure-flux integral in terms of $B_R$ alone, and the Lean estimate bounds it in terms of $B_R$ and also $A_R$. The source reports that the two are not directly comparable, since equation (10.17) bounds $B_R$ using $A_R$ but cannot be used to eliminate $A_R$ from the Lean estimate to recover (10.19). On that report neither content entails the other, and the transition is classified as lateral. The quantifier on time also differs: the natural-language estimate is for almost every $t$ and the Lean estimate for every $t$ (Chapter~\ref{ch:pressure}), which together with the extra dependence on $A_R$ is what makes the pair incomparable and not merely ordered. The classification is source-supported and provisional: the incomparability is the source's assessment of the two bounds, and no countermodel or analytic argument establishing both non-entailments is supplied in this monograph. The same holds for the converse non-entailment in the derivative case above. The example matters because it shows that the typing is not a ranking by strength. A discrepancy between a natural-language lemma and its formalization need not be a loss of content in the sense of weakening, and the framework represents the difference without flattening it into a generic mistranslation. As in the derivative case, no conclusion about the correctness of the full natural-language proof is drawn, in keeping with the source's own disclaimer: the Lean results do not refute the natural-language estimates, and they fail to certify them.

\section{Propagation through dependency structures}

A formal development is a graph of declarations, each depending on definitions and lemmas that precede it. Correspondence obligations attach to individual declarations: the definition of a quantity corresponds, or does not, to the notion in the source. The question of RV-013 is how such obligations propagate along the dependency graph.

Say that a declaration $x$ \emph{semantically depends} on declarations $y_1, \dots, y_k$ when there are maps $g_x$ on denotations and $g_x^{\flat}$ on formal counterparts with $\llbracket x \rrbracket_{\mathcal{I}} = g_x(\llbracket y_1 \rrbracket_{\mathcal{I}}, \dots, \llbracket y_k \rrbracket_{\mathcal{I}})$, the formal counterpart of $x$ being obtained by applying $g_x^{\flat}$ to the formal counterparts of the $y_j$. A \emph{congruence certificate} for $x$ is an accepted certificate establishing that this construction preserves correspondence under the intended interpretation: whenever the formal counterpart of each $y_j$ corresponds to $y_j$, the formal counterpart of $x$ corresponds to $x$.

\begin{proposition}[RV-013, conditional propagation of correspondence]
\label{prop:rv013}
Let $x$ semantically depend on $y_1, \dots, y_k$. (i) Suppose the correspondence obligation of each $y_j$ is warranted-discharged, and a congruence obligation $c_x$, of dependency type with the content that correspondence of all the $y_j$ implies correspondence of $x$, is warranted-discharged on the strength of an accepted congruence certificate. Then the correspondence obligation of $x$ may be discharged by a certificate whose dependencies are $c_x$ and the correspondence obligations of the $y_j$, and that discharge is currently warranted as long as $c_x$ and each $y_j$ remain warranted-discharged and no defeater targets the new certificate. If the new certificate is the only one for $x$, the obligation of $x$ is warranted-discharged exactly then. (ii) Suppose some $y_j$ has a correspondence obligation in the residue, and no invariance or bypass certificate establishes correspondence of $x$ independently of $y_j$. Then this dependency-based route does not justify discharge of the correspondence obligation of $x$.
\end{proposition}

\begin{proof}
(i) The congruence obligation $c_x$ has the content of a conditional: correspondence of every $y_j$ implies correspondence of $x$. Its antecedent is supplied by the discharged obligations, and the conclusion follows by modus ponens, the new discharge having $c_x$ and the obligations of the $y_j$ as its dependencies, which lie in the allowed set of obligations and authorized assumptions. By Definitions~\ref{def:warranted} and~\ref{def:status} its warrant is the conjunction of theirs together with the absence of a defeater on the certificate itself, and it is lost if any of them loses warrant. (ii) The route consists of the congruence certificate together with the correspondence of all the $y_j$. If one antecedent is not established, the conditional cannot be detached, and no other premise of the route bears on $x$.
\end{proof}

Discharging the correspondence obligations of all dependencies is not sufficient by itself, and the congruence certificate is therefore an explicit hypothesis of (i). A declaration constructed from faithful parts may fail to be faithful if the construction does not preserve the relevant semantic relation, for example if it applies a faithful definition under a quantifier whose range differs from the source's. The second clause is deliberately weaker than the statement that $x$ cannot be faithful. An unresolved dependency may turn out to be irrelevant to $x$, either because the denotation of $x$ does not in fact vary with it or because $x$ can be shown to correspond by a route that does not use it. That irrelevance is a further claim, and it is recorded by an invariance or bypass certificate, which is itself an obligation. The proposition says that the dependency graph cannot be used as evidence of correspondence when a dependency is unresolved, and that the residue of $x$ includes, in the sense of a recorded route, the residue of the dependencies it semantically uses. A single silent default in a definition, as in the hidden-existence family of the preceding chapter, thereby enters the residue of every statement that uses the defined quantity unless an invariance certificate removes it.

\section{Optional information and the link to admissible degradation}

The earlier work on admissible degradation concerns continuation under loss: a process may drop information without invalidating what it continues to do, provided what is dropped is not what justifies the continuation. The preservation condition provides the formal counterpart for obligations. A transformation may drop content that no originating obligation requires, which in the present vocabulary is optional content, and such a transformation can still be conservative or strengthening at the obligations that matter, since preservation is assessed against the contents of those obligations. It may not silently discard an obligation that justifies the authority of the continued artifact, which is the case in which a weakening or lateral transition is mislabeled as preserving. The distinction between permitted and forbidden loss is thereby reduced to the entailment condition, with the set of originating obligations playing the role of the authority-bearing content.

\section{What remains}

The principal mathematical work that remains concerns the certificates themselves. The theorems above are conditional on the soundness of the certificates that witness entailment and discharge, and the theory does not say when a certificate justifies what it claims. Three questions are open. The first is how to certify an entailment between natural-language contents and formal contents, where the consequence relation is relative to an interpretation that is itself informal. The second is how to certify that a refinement of a specification, such as the introduction of a more explicit definition, preserves the contents that depended on the original. The third is how the interpretation is itself allowed to change along a chain, which the present formalism holds fixed and which the audit section of the appendix records as a limitation. These questions are the subject of the constructive proposals of Part~VII, and the Lean prototype described in Chapter~\ref{ch:system} is designed to test whether recording obligations and transition certificates improves the reliability of certification claims. The criterion for that test is narrow. It is whether known unresolved obligations are preserved or detected and are not reported as settled.
